% Formal Grammar
\section{Formal Grammar of Indus Inscriptions}

\subsection{Context-Free Grammar Approximation}

While the Indus language is regular, we can describe it with a restricted CFG:

\subsubsection{Production Rules}
\begin{align*}
S &\rightarrow \text{INST} \mid \text{INST}~M \\
M &\rightarrow \text{COUNT} \mid \text{REL} \mid \text{QUAL} \mid \text{CLOSE} \\
M &\rightarrow \text{COUNT}~M \mid \text{REL}~M \mid \text{QUAL}~M \\
M &\rightarrow \text{CLOSE}
\end{align*}

\subsection{Regular Grammar}

More precisely, as a right-linear grammar:

\begin{align*}
S &\rightarrow \text{inst}~A \mid \text{inst} \\
A &\rightarrow \text{count}~A \mid \text{count}~B \mid \text{count} \\
A &\rightarrow \text{rel}~A \mid \text{qual}~A \\
B &\rightarrow \text{close}
\end{align*}

This generates exactly the regular language of valid inscriptions.

\subsection{Extended Backus-Naur Form (EBNF)}

Using modern notation:
\begin{verbatim}
inscription = instantiation [operations] [closure];
operations = {counting | relation | qualification};
counting = stroke+;
stroke = "|";
closure = terminal_sign;
\end{verbatim}

\subsection{Constraint Grammar}

Additional constraints beyond context-free:

\subsubsection{Type Constraints}
\begin{itemize}
    \item $\text{COUNT}$ must follow compatible $\text{INST}$
    \item $\text{REL}$ must have appropriate left and right arguments
    \item $\text{QUAL}$ must modify preceding entity
\end{itemize}

\subsubsection{Cardinality Constraints}
\begin{itemize}
    \item At most one $\text{INST}$ per inscription
    \item At most one $\text{CLOSE}$ per inscription
    \item Unbounded $\text{COUNT}$ (but practically limited)
\end{itemize}

\subsection{Parsing Algorithm}

A simple shift-reduce parser suffices:

\begin{verbatim}
function parse(inscription):
    state = INITIAL
    for sign in inscription:
        if can_shift(state, sign):
            state = shift(state, sign)
        else:
            return ERROR
    return state in ACCEPTING_STATES
\end{verbatim}

\subsection{Ambiguity}

The grammar is unambiguous: each valid inscription has exactly one parse tree (proof by determinism of transition function).

\subsection{Language Recognition}

The grammar generates exactly the language:
\[
L_{\text{Indus}} = \{w \in \Sigma^* \mid w \text{ is a valid inscription}\}
\]

This language is:
\begin{itemize}
    \item Regular
    \item Decidable
    \item Not context-free (CFL properties unnecessary)
    \item Finite (bounded by sign inventory and length constraints)
\end{itemize}
